\section{Documentation de l'existant}

% ----------------------------------------------------------------
% 📌 À quoi sert cette section ?
% ----------------------------------------------------------------
% Cette partie sert à prouver que vous avez exploré les solutions déjà disponibles
% (outils, logiciels, méthodes, articles) en lien avec le problème que vous abordez.
% Elle permet de montrer que votre projet est justifié par un manque, une faiblesse
% ou une opportunité d'amélioration dans ce qui existe déjà.

% ----------------------------------------------------------------
% 🧭 Comment la remplir ?
% ----------------------------------------------------------------
% ➤ Listez les outils ou approches existants en lien avec votre projet.
% ➤ Décrivez leurs fonctionnalités principales.
% ➤ Soulignez leurs avantages ET leurs limites dans le cadre de votre problématique.
% ➤ Comparez-les entre eux si pertinent.
% ➤ Concluez par l’intérêt de votre projet comme complément, alternative ou amélioration.

% Vous pouvez aussi ajouter un tableau comparatif ou un schéma d’analyse.

% ----------------------------------------------------------------
% 🧾 Exemple rédigé :
% ----------------------------------------------------------------
% Plusieurs plateformes de test en ligne existent pour manipuler du code SVG :
% - \textbf{CodePen} : puissant et flexible, mais non spécialisé SVG
% - \textbf{SVGPlayground} : permet de tester du SVG, mais sans objectif éducatif
% - \textbf{JSFiddle} : généraliste, sans guidage pédagogique
%
% Ces outils ont en commun leur accessibilité, mais manquent d’accompagnement pour des profils débutants.
% Aucun ne propose une approche gamifiée ou structurée autour de défis pédagogiques.

% ----------------------------------------------------------------
% 📊 Exemple de tableau comparatif
% ----------------------------------------------------------------
% \begin{table}[H]
% \centering
% \begin{tabular}{|l|c|c|c|}
% \hline
% Outil & Ciblé débutant & Pédagogique & Temps réel \\
% \hline
% CodePen & Non & Non & Oui \\
% SVGPlayground & Partiellement & Non & Oui \\
% Projet proposé & Oui & Oui & Oui \\
% \hline
% \end{tabular}
% \caption{Comparaison des solutions existantes}
% \end{table}

% ----------------------------------------------------------------
% ✍️ À vous de jouer : complétez ci-dessous
% ----------------------------------------------------------------

% (Début de réponse)
À ce jour, plusieurs outils existent dans le domaine, parmi lesquels...

